\clearpage
\fancyhead[L]{\small LAPLACE. EXPOSITION DU SYSTÈME DU MONDE.}
\SourceLeaf{505}
\begin{center}\small Göttingische gelehrte Anzeigen. No. 119. pp. 1185..1192. 23 July 1808.\\[3pt]\rule{.72\linewidth}{.3pt}\end{center}

\emph{Exposition du système du monde, par} M. \textsc{Laplace}, \emph{chancelier du Senat-Conservateur etc. Troisième édition, revue et augmentée par l'auteur. Chez Courcier,} 1808. 405 pages in quarto. With a portrait of the author.

The subject, plan, approach and value of this work, a classic of its kind whose earlier editions are in everyone’s hands, are too well known for us to dwell on them now. A detailed review of its contents, based on the German translation of the first edition, can already be found in issues 97 and 131 of these pages for 1798. The second edition, which appeared in 1799, had undergone no significant changes. We shall therefore confine ourselves here to the principal changes and additions that distinguish this third edition from the earlier ones. The increase in the number of pages already shows that the additions must be substantial, since the second edition had only 351 pages. They consist partly of a few larger additions and partly of a great many smaller ones. These make the treatment of the subjects more complete and their exposition more coherent, clearer, more fruitful and more interesting. A particular merit of this work is that it gives the latest and best-established numerical results of astronomy, insofar as such results belong in a general outline. In this respect it can all the more readily be regarded as an authority, since many of these determinations were established or prompted by the author’s own profound investigations. We believe that we shall do some readers a welcome service by picking out in this review the most important of those determinations that have been further improved since the 2nd edition.

The \emph{first book}, which deals with the \emph{apparent} motions of the celestial bodies, has gained two chapters. \emph{The tenth} is on the telescopic planets, Ceres, Pallas, Juno and Vesta, and \emph{the eleventh} is on the motion of the planets around the Sun. The second chapter of the second book, which had the same heading, has been omitted. The grounds for the true arrangement of the world do indeed stand out more clearly once it has been shown that the planets’ motions appear to us exactly as if they took place in epicycles whose centres always coincided with the Sun and moved with it around the Earth. For Venus and Mercury this can be proved from their phases together with the successive changes in their apparent diameters. For Jupiter and Saturn it can be proved from the eclipses of the satellites and the disappearances of the ring. These conclusions drawn from the facts therefore rightly belong in the first section, since they are independent of the choice of world system. Before the invention of telescopes, of course, these facts were unknown. In his system \textsc{Ptolemy} could therefore give only the ratio of the radii of the circles and epicycles for each planet and had to leave their absolute sizes undetermined.

We note the following improved numerical results from this book, referring to the pages of the second edition. The new decimal divisions are used throughout, as scarcely needs mentioning. On p. 6, the obliquity of the ecliptic for 1801 is $26^{\circ}\,07315$. Its secular decrease for the present century (p. 11) is $160''85$. On p. 16, the sidereal day is 0.997269672 of the mean solar day. On p. 17, the tropical year is now 365.242264 days. On p. 20, the sidereal revolution of the Moon takes 27.3216610716 days at the beginning of the 19th century. The sidereal revolution of the perigee takes 3232.58075
\SourceLeaf{506} days and that of the node 6793.42118 days. At the epoch for 1801 (midnight before 1 January in Paris), the mean longitude of the Moon is $124^{\circ}01299$, the perigee $295^{\circ}66824$ and the ascending node $17^{\circ}6933$. The eccentricity of the Moon’s orbit is 0.0548553, its inclination $5^{\circ}7222$, and the synodic revolution of the Moon takes 29.58058817896 days. We shall bring together the new determinations concerning the planets below. On p. 51, the major axis of the nutation ellipse is $59''56$ and the minor axis is $111''30$ in units of the parallel of latitude. On p. 55, the northern degree at mean latitude $73^{\circ}7$, which measured 100696 metres according to the French astronomers, comes to only 100316.1 metres according to the new Swedish arc measurement. A repeated measurement might also alter considerably the length of the degree measured by \textsc{Lacaille} at the Cape. It seems that perhaps rather too much weight has been placed on it as evidence of the irregularity of the Earth’s figure. To the account of the principal methods for determining longitude, \textsc{Laplace} adds the wish that all nations might agree to reckon longitude from some particularly distinctive natural landmark instead of from their chief national observatory. The peak of Tenerife would be especially well suited to this. On p. 81, the ratio of the specific weight of atmospheric air to that of mercury, at a barometer reading of 0.76 and the temperature of melting ice, is 1 to 10477.9. The whole section on measuring heights with the barometer and on the refraction of light is treated in greater detail.

In the \emph{second book}, on the true motions of the celestial bodies, we find few changes. In chapter VI (now V) the estimate of the probability that two observed comets with nearly identical elements are one and the same has been omitted, and in its place some conjectures about the action of the Sun’s heat on comets have been ventured. The table of the defining quantities for all the planets, on the other hand, has been completely revised. We therefore reproduce it here in full on account of its importance. We have already mentioned that \textsc{Laplace} now reckons the epoch from midnight between 31 December and 1 January.

\begin{center} \Needspace{140pt}
{\small\setlength{\tabcolsep}{6pt}
\begin{tabular}{l|r|r}
 & \shortstack{Mean longitude 1801\\Paris meridian} & \shortstack{Sidereal period\\in days} \\
\hline
Mercury & $182^{\circ}15647$ & $87.96925804$ \\
Venus & $11.93672$ & $224.70082399$ \\
Earth & $111.28179$ & $365.25638350$ \\
Mars & $71.24145$ & $686.9796186$ \\
Jupiter & $124.67781$ & $4332.5963076$ \\
Saturn & $150.38010$ & $10758.9698400$ \\
Uranus & $197.54244$ & $30688.7126872$ \\
\end{tabular}} \end{center}

\begin{center} \Needspace{140pt}
{\small\setlength{\tabcolsep}{6pt}
\begin{tabular}{l|r|r|r}
 & \shortstack{Semimajor axis} & \shortstack{Eccentricity\\1801} & \shortstack{Secular change} \\
\hline
Mercury & $0.3870981$ & $0.20551494$ & $+ \text{c}.000003867$ \\
Venus & $0.7233323$ & $0.00685298$ & $- 0.000062711$ \\
Earth & $1.0000000$ & $0.01685318$ & $- 0.00\text{c}041632$ \\
Mars & $1.5236935$ & $0.09313400$ & $+ 0.000090176$ \\
Jupiter & $5\;2027911$ & $0.04817840$ & $+ 0.000159350$ \\
Saturn & $9.5387705$ & $0.05616830$ & $- 0.000312402$ \\
Uranus & $19.1833050$ & $0.04667030$ & $- 0.000025072$ \\
\end{tabular}} \end{center}

\SourceLeaf{507}
\begin{center} \Needspace{140pt}
{\small\setlength{\tabcolsep}{6pt}
\begin{tabular}{l|r|r}
 & \shortstack{Perihelion 1801} & \shortstack{Sidereal\\secular motion} \\
\hline
Mercury & $82^{\circ}6256$ & $+ 1801''10$ \\
Venus & $142.9077$ & $- 826.63$ \\
Earth & $110.5571$ & $+ 3641.40$ \\
Mars & $369.3407$ & $+ 4884.05$ \\
Jupiter & $12.3812$ & $+ 2048.95$ \\
Saturn & $99.0549$ & $+ 5978.60$ \\
Uranus & $185.9574$ & $+ 738.69$ \\
\end{tabular}} \end{center}

\begin{center} \Needspace{140pt}
{\small\setlength{\tabcolsep}{6pt}
\begin{tabular}{l|r|r}
 & \shortstack{Orbital inclination\\1801} & \shortstack{Secular change} \\
\hline
Mercury & $7^{\circ}78058$ & $+ 56''12$ \\
Venus & $3.76936$ & $- 14.05$ \\
Mars & $2.05663$ & $- 0.47$ \\
Jupiter & $1.46034$ & $- 69.78$ \\
Saturn & $2.77102$ & $- 47.88$ \\
Uranus & $0.85990$ & $+ 9.67$ \\
\end{tabular}} \end{center}

\begin{center} \Needspace{140pt}
{\small\setlength{\tabcolsep}{6pt}
\begin{tabular}{l|r|r}
 & \shortstack{Ascending node\\1801} & \shortstack{Sidereal\\secular motion} \\
\hline
Mercury & $51^{\circ}0651$ & $- 2414''41$ \\
Venus & $83.1972$ & $- 5770.99$ \\
Mars & $53.3605$ & $- 7186.65$ \\
Jupiter & $109.3624$ & $- 4869.04$ \\
Saturn & $124.3662$ & $- 6995.25$ \\
Uranus & $10.9488$ & $- 11104.81$ \\
\end{tabular}} \end{center}

Of the elements given here for the four new planets, the sets for Ceres, Pallas and Juno are the 11th, 9th and 6th by \textsc{Gauss}, respectively. For Vesta they are the ones Mr \textsc{Burkhardt} gave in the Connoissance des tems for 1809, for which only a few observations had been used at the time. We therefore leave them out here, since more accurate ones have long been known in Germany (Monatl. Corresp. Febr. 1808; Gött. gel. Anz. 1808 St. 14. 40. 107).

The \emph{third book}, on the laws of motion, has likewise received only a few small additions. They concern the expression of the first principles of dynamics in equations, the principle of least action, the concept of \emph{mass} and the two kinds of equilibrium.

In the \emph{fourth book}, on the theory of universal gravitation, the chapter on the perturbations of the planets’ elliptical motions has become the second chapter. Two new chapters have been added, one on the satellites of Saturn and Uranus and one on the attraction between the smallest parts of bodies that is perceptible only at imperceptible distances (attraction moléculaire). This last chapter gives an overview of the theory that \textsc{Laplace} developed in two supplements to his Mécanique céleste (De l'action capillaire, and Supplément à la théorie de l'action capillaire). He applied it with such surprisingly happy results to the mathematical explanation of the refraction of light, phenomena in capillary tubes, the attraction and repulsion of small bodies floating on a liquid, the adhesion of a disk to a liquid, the flotation of small solid bodies in a liquid of lower specific weight, and so on. These discoveries mark an epoch in physics. They open the way to a science that will be to nature in its smallest parts what the theory of gravitation is to nature on the large scale. We must reserve a fuller account of this for another occasion.

\SourceLeaf{508}
Among the smaller changes and additions in this section we also note the following. The great equation of Saturn reaches $9111''41$, and its period is $921\frac{1}{2}$ years. The great equation of Jupiter is $3720''36$. The masses of the planets are to the mass of the Sun as 1 to 2025810 for Mercury, 356632 for Venus, 337086 for the Earth, 2546320 for Mars, 1067.09 for Jupiter, 3534.08 for Saturn and 19504 for Uranus. The apparent diameter of Jupiter at distance 1 is $599''151$. The weights of a body at the equator of the Earth, at that of Jupiter and at that of the Sun are in the ratio 1000, 2566, 27933. Several additions have been inserted into the theory of the Moon. They concern the inequalities that depend on the flattening of the Earth and on the parallax of the Sun, the secular inequalities in the motion of the node and the apsides, and the equation discovered a few years ago whose period is 184 years. --- The section on Jupiter’s satellites has also been expanded in several respects.

The \emph{fifth book} contains the history of astronomy. Here too there are only outlines, but drawn by such a hand they hold great interest for the expert and the amateur alike. “An account of the advances of the most sublime of the natural sciences,” says \textsc{Laplace}, “advances that continue even amid the revolutions of empires, may console us for the misfortunes whose stories fill the annals of all peoples.” This part of the work, too, has been enriched by several additions and interesting reflections. Among them are the remarks on the traces of astronomical knowledge found among the Indigenous peoples of Mexico and Peru. The former knew the length of the tropical year very precisely. Their method of intercalation assumes a year of $365\frac{25}{104}$ days, or 365 days, 5 hours, 46 minutes and 9 seconds. It is quite remarkable that they did not know the division of time into weeks, which is found among all the peoples of the Old World, and that in its place they had a period of five days.

There are also six historical notes at the end of the work. The first concerns a couple of Chinese observations by \textsc{Tcheou-Kong}, according to which around the year 1100 before our era the obliquity of the ecliptic was $26^{\circ}5563$ and the right ascension of the star $\varepsilon$ in Aquarius was $297^{\circ}8096$. According to \textsc{Laplace}’s formulas, the former should have been $26^{\circ}5161$ and the latter $298^{\circ}7265$ in that year. To bring the latter numbers into full agreement, one would only have to go back another 54 years. The second concerns an account preserved for us by \textsc{Geminus} of what the Chaldeans knew about the Moon’s motion. The third concerns the observation of the solstice at Marseille by \textsc{Pytheas}. The fourth compares the values given by \textsc{Hipparchus} for the Moon’s motion with respect to the Sun, the perigee and the node with the determinations that follow from \textsc{Laplace}’s theory of secular inequalities. The fifth compares the greatest equation of the centre of the Sun, the position of aphelion, the length of the year and the obliquity of the ecliptic as determined by the Arabs with the latest values. Finally the sixth sets the determinations of the obliquity of the ecliptic by \textsc{Tcheou-Kong} in 1100 BC, \textsc{Pytheas} in 350 BC, \textsc{Ibn Junis} in 1000, \textsc{Cocheou-Kong} in 1280 and \textsc{Ulug-Beigh} in 1437, and the latest ones from 1801, alongside the results of the theory.
